% ============================================================================
\clearpage
\section*{v4.9.0-r77 Source-Only: High Precision, Formal Integrability at the Germ and the Covariance of the OP5 Comparison System}
\addcontentsline{toc}{section}{v4.9.0-r77 Source-Only: High Precision, Formal Integrability at the Germ and the Covariance of the OP5 Comparison System}
\label{sec:v490-r77-high-precision-covariance}
% ============================================================================

\paragraph{Purpose.}
Layer v4.9.0-r74 left two questions open (\S\ref{sec:v490-r74-precision-limit}).  The first is whether the extended symbol
\(g''\) is involutive first at order 8.  The second is whether the system has further hidden conditions on 7- and 8-jets;
Spencer allows up to four and three.  Both were blocked by double precision.  On 3 and 4 November 2026 the author asked
for three things:
\begin{itemize}[leftmargin=2em]
\item to repeat the whole chain in high precision;
\item to run the decisive eliminations on parallel GitHub runners;
\item to look for links of the result to physics.
\end{itemize}
This layer records one high-precision exploration and six pre-registered tests (H59 to H64).  Three tests are confirmed
and three refuted.  Everything concerns the comparison model of \S\ref{sec:v410-state-mediator} only.  Earlier layers are
not rewritten.

\noindent\textbf{Source-only status.}
This v4.9.0-r77 layer is source-only.  It does not build a PDF, does not assign a new DOI,
and does not authorize a public release.  The v4.9.0 Version DOI
\href{https://doi.org/10.5281/zenodo.23534954}{10.5281/zenodo.23534954} and the Concept DOI
\href{https://doi.org/10.5281/zenodo.19210728}{10.5281/zenodo.19210728} are unchanged by it.

\subsection*{v4.9.0-r77 results}
\label{subsec:v490-r77-results}

\paragraph{Record v4.9.0-r77.1 (exploration: the four conditions and the Cartan test in high precision).}
Not pre-registered and not blind: the values of r74.1 were known.
\begin{itemize}[leftmargin=2em]
\item \emph{Linearization.}  The linearization was rebuilt pointwise in 160 digits.  The deviation from the earlier
high-precision data is \(2\times10^{-21}\) in degree 1 and \(2\times10^{-19}\) in degree 2; these earlier data set the old
precision limit.
\item \emph{Eliminations.}  \(c_8\) and \(c_7\) were eliminated in ball arithmetic with 70 digits.  The Newton residuals are
\(10^{-66}\) and \(9\times10^{-52}\).
\item \emph{Fibre of 6-jets.}  The Gram matrix of the derived conditions on the 501-dimensional fibre has 24 genuine
eigenvalues, the smallest \(5.1\times10^{-13}\), then \(1.9\times10^{-71}\): a gap of \(10^{58}\).  The fibre therefore has
dimension exactly \(477\).  It consists of 20 \(\Omega\)-prolongations and exactly four new conditions.
\item \emph{Cartan test.}  On this basis \(\dim g''_p=477,594,723,865\) for \(p=6,\dots,9\).  The characters at \(p=8\) are
\(723,130,12,0\) in all three frames.  The symbol is not involutive at \(p=6,7\) and is involutive at \(p=8\), with rank
gaps \(\ge10^{8}\).
\item \emph{A withdrawn local finding.}  An intermediate local finding (only 48 instead of 112 identities at order 4) came
from a conversion with 15 digits and was withdrawn before this layer.
\end{itemize}

\paragraph{Theoremlet v4.9.0-r77.2 (H59 confirmed: no further hidden conditions on 7- and 8-jets).}
Pre-registered, not blind (the symbol dimensions were known).  Before the decisive runs, a dated addendum relaxed the
residual bound from \(10^{-40}\) to \(10^{-30}\), after a function test at \(q=4\).
\begin{itemize}[leftmargin=2em]
\item \emph{Setup.}  The jet through degree \(q-1\) is fixed and the equations go through degree \(q+1\).  \(c_{q+2}\) and
\(c_{q+1}\) are eliminated in 70 digits on parallel runners (16 GB each).
\item \emph{Result.}  The fibres are \(F_6=477\) (positive control), \(F_7=637-43=594\) and \(F_8=789-66=723\).  The Gram gaps
are \(1.6\times10^{58}\), \(2.8\times10^{58}\) and \(3.9\times10^{56}\); all residuals are \(\le1.1\times10^{-51}\).
\item \emph{Reading.}  The Spencer bounds of r74.2 are not attained.  Together with the involutivity of \(g''_8\), the
completed system is formally integrable at the germ by Goldschmidt's criterion, provided two elimination levels catch all
obstructions as at order 6.  The general formal solution depends on 12 functions of three variables.
\end{itemize}

\paragraph{Theoremlet v4.9.0-r77.3 (H60 confirmed: six polarizations per null direction).}
Pre-registered.  With \(s_3=12\) and a cone of degree 2, the multiplicity on the cone is 6.  If the symbol module is reduced
there, this gives six plane waves per null direction.  Over 200 null directions at each of \(p=6,\dots,9\), the polarization
space \(\{u:\xi^p\otimes u\in g''_p\}\) has dimension 6 every time, with gaps \(\ge3\times10^7\).  Off the cone the dimension
is 0.  The twelve Cauchy data are therefore six real massless modes.  The deficit of plane waves seen in r74.4 comes from the
planes \(\Pi_0,\Pi_1\) and from isolated directions, not from the generic cone.

\paragraph{Theoremlet v4.9.0-r77.4 (H61 refuted: no Weyl type in the naive spin assignment).}
Pre-registered.  The test reads \(\Psi\) as three Dirac spinors in the Dirac basis, with spin \(I_3\otimes\Sigma\), chirality
\(I_3\otimes\gamma_5\) and complex structure \(J\).
\begin{itemize}[leftmargin=2em]
\item The six-dimensional polarization space is invariant under none of \(J\), \(\gamma_5\) or the helicity: the median
residuals are \(1.0\), \(0.99\) and \(0.5\).
\item Post hoc (exploration): \(P+JP\) has dimension 12 but is not the kernel of the massless Dirac symbol, so a Majorana
reading also fails.  The modes spread over all three blocks of \(\Psi\), most of them over the third block, which does not
enter the coframe.
\end{itemize}

\paragraph{Theoremlet v4.9.0-r77.5 (H62 refuted: the symbol is invariant only under scaling).}
Pre-registered.  The stabilizer
\(\{(A,B)\in\mathfrak{gl}(4)\times\mathfrak{gl}(24):X_{A,B}g''_p\subset g''_p\}\) has dimension 50 at \(p=6,7,8\), with gaps
\(\ge7\times10^8\).  Its spacetime part is only \(A\propto1\), and its intersection with \(\mathfrak{so}(Q)\) is 0, not 6.  The
internal part (dimension 49) is closed under commutators.  Its Killing form has signature \((3,1,45)\), so the internal part
is mostly nilpotent; no compact semisimple part is visible (exploration).  A first run that formed the Gram matrix as a
difference cancelled digits and is not counted.  An addendum places the result: invariance at a fixed background is expected
to fail, because every symmetry must preserve the characteristic set, and \(J(\xi)\) loses rank on \(\Pi_0\cup\Pi_1\).  These
planes are not \(q\)-orthogonal, are not spanned by bilinears of \(\Psi\) (r71, \texttt{k122}) and arise in the \(G\) part
(\texttt{k126}).

\paragraph{Theoremlet v4.9.0-r77.6 (H63 refuted: already \(F\) is not Lorentz symmetric at a fixed background).}
Pre-registered.  For the blocks of \(J\) the stabilizer
\(\{(A,B,C):\sum_\mu A_{\mu\nu}X_\mu+X_\nu B-CX_\nu=0\}\) has spacetime parts of dimension 7 for \(F\) and 1 for \(F+G\) and
\(F+G+dg\).  In all three cases the intersection with \(\mathfrak{so}(Q)\) is 0.  Part (a), predicting 6 for \(F\), fails;
part (b), predicting \(\le1\) when \(G\) is included, holds.  The six non-scaling directions of \(F\) lie near
\(\mathfrak{so}(Q)\) (cosines 0.74 to 0.9998) but form no \(\mathfrak{so}(Q')\) for any form \(Q'\) (exploration).  The germ
\(\Psi\) enters the principal symbol beyond a tetrad, through the coframe built from \(\Psi\).  The symbol is then invariant
only under the stabilizer of \(\Psi\), which is trivial for a generic spinor.

\paragraph{Theoremlet v4.9.0-r77.7 (H64 confirmed: \(R\) is covariant).}
Pre-registered.  Run 1 was invalid: its positive control was \(2.1\times10^{-7}\), above the bound \(10^{-9}\), because of an
unregularized pseudoinverse of the rank-31 symbol.  A dated addendum preceded run 2, which is not blind.  The construction
\(R(\Psi,\partial\Psi;A,\partial A,h,\partial h)\) was rebuilt after three transformations: a random linear change of
coordinates in \(GL(4)\), a random internal Lorentz transformation of \(\Psi\), and both together.  In each case the jets and
the background transform tensorially.
\begin{itemize}[leftmargin=2em]
\item With a free row transformation \(C\), the relative residual \(\min_C\|J'L-CJ\|/\|J'L\|\) is
\(\le4.1\times10^{-10}\) for \(F\), \(G\) and \(dg\).
\item The positive control gives \(\le1.1\times10^{-11}\).
\item The negative control, with the background held fixed, gives \(0.26\) for \(F\) and \(G\).
\end{itemize}
At symbol level, \(R\) is therefore coordinate covariant and internally Lorentz covariant.  The breaking found in H62 and H63
comes entirely from the background \((\Psi\text{ at the germ},A,h)\).

\paragraph{Proof and measurement sources.}
\begin{itemize}[leftmargin=2em]
\item r77.1: \texttt{k153b\_hp\_linear.py}, \texttt{k154\_hp\_tools.py}, \texttt{k157\_hp\_bedingungen.py},
\texttt{k158\_cartan\_hp.py}.
\item r77.2: \texttt{k159\_hp\_faser.py}, with the workflow \texttt{hp-faser.yml} and the data under \texttt{ci/hp/}.
\item r77.3: \texttt{k160\_polarisationen.py}.
\item r77.4: \texttt{k161\_spin.py}, \texttt{k161b\_majorana.py}, \texttt{k161c\_bloecke.py}.
\item r77.5: \texttt{k162\_symmetrien.py}, \texttt{k162b\_innere.py}.
\item r77.6: \texttt{k163\_lorentz\_F.py}, \texttt{k163b\_eigene\_metrik.py}.
\item r77.7: \texttt{k164\_kovarianz.py}.
\end{itemize}
The pre-registration notes \texttt{Vorab\_\allowbreak Notiz\_\allowbreak H59} to \texttt{H64} and the result notes lie in
\texttt{notes/\allowbreak research/\allowbreak predictions/}.  \(\square\)

\paragraph{Consequence.}
The extended comparison system is formally integrable at the germ, with six real massless modes.  Its construction is
covariant; Lorentz symmetry is broken only by the background.  The physical questions therefore move to the background:
\begin{itemize}[leftmargin=2em]
\item what the fields \(A\) and \(h\) are;
\item whether the background solves covariant equations;
\item whether \(\Omega\), \(\theta\), the gauge \(\Psi^\dagger K_2\Psi=S\) and the four conditions are covariant;
\item which Lorentz action, if any, acts on the six modes.
\end{itemize}
OP5 is not closed.

\subsection*{v4.9.0-r77 status in the three frames}
\label{subsec:v490-r77-frames}

\begin{itemize}[leftmargin=2em]
\item All statements concern the fixed OP5 comparison model in one chart, with one coframe reconstructed from \(\Psi\).  No
frame map M, T or P acts on them.  The Lorentz transformations of r77.5 to r77.7 act inside the model and are not the photon
frame P.
\end{itemize}
No statement of this layer concerns colour or charge.  The internal \(3\otimes\) structure of \(\Psi\) is not read as
generations.

\subsection*{v4.9.0-r77 register entries}
\label{subsec:v490-r77-register}

\begin{center}
\small
\begin{tabular}{@{}>{\raggedright\arraybackslash}p{0.34\textwidth}>{\raggedright\arraybackslash}p{0.58\textwidth}@{}}
\toprule
statement & status after v4.9.0-r77\\
\midrule
four hidden conditions & exactly four (gap \(10^{58}\), high precision, exploration)\\
extended symbol & involutive first at \(p=8\), characters \(723,130,12,0\) (exploration)\\
further hidden conditions & none on 7- and 8-jets (H59 confirmed); formally integrable at the germ\\
massless modes & six real polarizations per null direction (H60 confirmed)\\
spin of the modes & not Weyl or Majorana in the naive assignment (H61 refuted); undetermined\\
symmetry at fixed background & only scaling (H62, H63 refuted); structurally expected\\
covariance of \(R\) & coordinate and internal Lorentz covariant at symbol level (H64 confirmed, run 2 not blind)\\
next step & covariance of \(\Omega\), \(\theta\), the gauge and the four conditions; origin of the background \(A,h\)\\
\bottomrule
\end{tabular}
\end{center}

\subsection*{v4.9.0-r77 ledger and guard update}
\label{subsec:v490-r77-ledger-guard}

\[
R_{\rm chk}^{v490\text{-}r77}=(1,1,0)
\]
for (formal integrability at the germ and covariance of \(R\) established; H61 to H63 recorded as refuted; OP5 not closed).

The layer spends no guard:
\[
R_{\rm glob}^{(v490\text{-}r77,+)}=(0,1,1,1,1),
\qquad
\Gamma_{\rm guard}^{v490\text{-}r77}=(1,1,1,1,1),
\qquad
\beta_4=1/256\ \text{unspent}.
\]
It does not imply exact finite \(C5\) zero closure, exact all-readout terminality,
\(C6/T2\), a global photon-clock/RH theorem, or global all-readout closure.
